\documentclass[10pt,a4paper]{amsart}
\usepackage[margin=19mm]{geometry}
\usepackage{amsmath,amssymb,mathtools}
\usepackage[scr=rsfs,cal=euler]{mathalfa}
\usepackage{xfrac}
\usepackage[unicode,hidelinks,bookmarks=false]{hyperref}
\newcommand{\ie}{i.e.\ }
\newcommand{\cf}{cf.\ }
\newcommand{\resp}{respectively\ }
\newcommand{\Subsection}{\subsection}
\providecommand{\nobreakdash}{\nobreak\hbox{-}}
\setlength{\emergencystretch}{2em}
\multlinegap0pt
\usepackage{bm}
\usepackage{bbm}
\usepackage{dsfont}
\usepackage{xspace}
\usepackage{cancel}
\usepackage{mathtools}
\usepackage{amsthm}
\makeatletter
\@ifundefined{Thm}{\newtheorem{Thm}{Thm}}{}
\@ifundefined{Coro}{\newtheorem{Coro}[Thm]{Corollary}}{}
\@ifundefined{Prop}{\newtheorem{Prop}[Thm]{Proposition}}{}
\makeatother
\newcommand\Creg{\mathsf{C}_{\textup{reg}}}
\newcommand\creg{\mathsf{c}_{\textup{reg}}}
\newcommand{\deltareg}{\delta_{\textup{reg}}}
\title[A random walk approach to high-dimensional critical phenomen]{A random walk approach to high-dimensional critical phenomena}
\author{Hugo Duminil-Copin, Aman Markar, Romain Panis, Gordon Slade}
\date{Source: arXiv:2605.21438v2 (CC BY 4.0)}
\begin{document}
\maketitle

\noindent\emph{Source and licence.} A random walk approach to high-dimensional critical phenomena. Version arXiv:2605.21438v2 (CC BY 4.0), distributed under the Creative Commons Attribution 4.0 International licence, as recorded in the arXiv licence metadata for this exact version and shown on the abstract page https://arxiv.org/abs/2605.21438v2. Full rights notice: licenses/arxiv-2605.21438-CC-BY-4.0.txt.\par\smallskip

\noindent\textit{Editorial scope.} Counted unit: the Prop whose source label is \texttt{prop: regularity below beta(delta)} in the article (source0.tex lines 1451--1454, source label \texttt{prop: regularity below beta(delta)}), delivered together with the article's own complete original proof of it (source0.tex lines 2430--2602, one \texttt{proof} environment of 173 lines). No mathematics is written, repaired or re-derived: statement and proof are the article's own text line for line. The proof is detached from the statement in the source: the article states the result at lines 1451--1454 and prints its proof at lines 2430--2602, and the proof environment's own header names the statement, so the association is the article's own. Required context retained uncounted: eq:R\_J comparable to xi(0) (lines 418--423); coro:Miteration (lines 1493--1506); prop: rough bounds chi xi (lines 2042--2055); eq:chi0 (lines 2070--2072); cor: comparison chi xi when E small (lines 2155--2165); prop:Zbd (lines 2226--2243). Every definition, display and label of those lines is retained unchanged. Omitted from this package and not redistributed: 1--417, 424--1450, 1455--1492, 1507--2041, 2056--2069, 2073--2154, 2166--2225, 2244--2429, 2603--6292. Nothing inside a retained excerpt is omitted or altered: every retained body is delivered exactly as the article's lines are written, so no omission receipt of any kind accompanies this package. Citations: the retained excerpts carry no citation command at all, so no bibliography entry of the article is transcribed and no reference is redistributed. Reference adaptation: none was needed and none was made; every reference inside the delivered unit resolves inside it, so no unresolved reference or citation is produced. Printed slips kept literal: no printed word, symbol, hypothesis or number of the counted statement or of its proof was altered. Layout and class: the article's own class is \texttt{article} with packages including inputenc, fontenc, lmodern, dsfont, indentfirst, amsfonts, amsthm, amsmath, amssymb, amscd, mathrsfs, eucal, graphicx, epstopdf; the wrapper is the lane's pinned \texttt{amsart} editorial wrapper at 10pt on A4, which restates the theorem-like environments the retained text needs and adds the article's own macro definitions that the retained text needs (\texttt{\textbackslash Creg}, \texttt{\textbackslash creg}, \texttt{\textbackslash deltareg}), with extra wrapper packages (bm, bbm, dsfont, xspace, cancel, mathtools). The delivered numbering is the unit's own. Assets: no retained span contains a figure environment, a graphics-inclusion command, a PDF-page inclusion, a table or a TikZ picture; no image file is copied into this package, the package declares no asset dependency and no image file is redistributed with it. Licence: version arXiv:2605.21438v2 of this article is distributed under the Creative Commons Attribution 4.0 International licence, as recorded in the arXiv licence metadata for this exact version and shown on the abstract page https://arxiv.org/abs/2605.21438v2. A full rights notice is delivered at licenses/arxiv-2605.21438-CC-BY-4.0.txt. Characters: every retained excerpt is pure ASCII, so the delivered engine, XeLaTeX with its default Latin Modern text font, reports no missing character.
\begin{equation}
\label{eq:R_J comparable to xi(0)}
    c_0R_J \le \sigma_J,
    \qquad
    J_{x} \le c_0^{-1}R_J^{-d}.
\end{equation}

\begin{Coro}
\label{coro:Miteration}
Let $0\leq \beta'\leq \beta<\beta_c$ and let $s \in \mathbb R$.
If $Z_{\beta',\beta}M_{\beta'}(s\xi(\beta')/\xi(\beta))) < 1$
then
\begin{equation}
\label{eq:Mfirstbd}
    M_\beta(s)
    \le
    \frac{\chi(\beta')}{\chi(\beta)}
    \frac{M_{\beta'}(s\xi(\beta')/\xi(\beta))}{1-Z_{\beta',\beta}M_{\beta'}(s\xi(\beta')/\xi(\beta))}
    .
\end{equation}
\end{Coro}

\begin{Prop}
\label{prop: rough bounds chi xi}
For every $0\leq \beta'\leq \beta < \beta_c$,
\begin{equation}\label{eq:bounds chi}
	(\beta-\beta')(1-E(\beta))\leq \frac{1}{\chi(\beta')}-\frac{1}{\chi(\beta)}\leq \beta-\beta',
\end{equation}
and, if $E(\beta)<1$,
\begin{equation}\label{eq:bounds chi xi} \left(\frac{\chi(\beta')}{\chi(\beta)}\right)^{\frac{1+E(\beta)}{1-E(\beta)}}
    \leq \left(\frac{\xi(\beta')}{\xi(\beta)}\right)^2
    \leq \left(\frac{\chi(\beta')}{\chi(\beta)}\right)^{ 1-2E(\beta)}.
\end{equation}
In particular, if $E(\beta)\le \frac 12$ then
$\xi(\beta') \le \xi(\beta)$.
\end{Prop}

\begin{equation}\label{eq:chi0}
	\frac{1}{1-\beta(1-\delta)}\leq \chi(\beta)\leq \frac{1}{1-\beta}.
\end{equation}

\begin{Coro}
\label{cor: comparison chi xi when E small}	
Let $0\leq \beta'\leq \beta< \beta_c$.
If $E(\beta)\leq \frac{1}{4}\frac{\chi(\beta')}{\chi(\beta)} \vee \frac{1}{4}(\frac{\xi(\beta')}{\xi(\beta)})^2$ and also $E(\beta)\leq \tfrac{1}{2}$, then
\begin{equation}
\label{eq: comparison chi xi when E small}
	\frac{1}{2}\frac{\chi(\beta')}{\chi(\beta)}
    \leq
    \left(\frac{\xi(\beta')}{\xi(\beta)}\right)^2\leq 2\frac{\chi(\beta')}{\chi(\beta)}.
\end{equation}
\end{Coro}

\begin{Prop}[Bounds on $Z_{\beta',\beta}$]
\label{prop:Zbd}
Let $0\leq \beta'\leq \beta<\beta_c$. Then,
\begin{equation}
\label{eq: bounds on Z1}
	1-\frac{\chi(\beta')}{\chi(\beta)}\leq Z_{\beta',\beta}\leq \Big(1-\frac{\chi(\beta')}{\chi(\beta)}\Big)\frac{1}{1-1\wedge E(\beta)}.
\end{equation}
In particular, if $E(\beta)\leq \tfrac{1}{4}\tfrac{\chi(\beta')}{\chi(\beta)}$, then
\begin{equation}\label{eq: explicit bounds on Z1}
	1-\frac{\chi(\beta')}{\chi(\beta)}\leq Z_{\beta',\beta}\leq 1-\frac{1}{2}\frac{\chi(\beta')}{\chi(\beta)}.
\end{equation}
If we assume instead
that $E(\beta)\leq \tfrac{1}{8}(\tfrac{\xi(\beta')}{\xi(\beta)})^2 \wedge \tfrac{1}{2}$, then
\begin{equation}
\label{eq: bounds on Z in terms of n/N}
	1-2\left(\frac{\xi(\beta')}{\xi(\beta)}\right)^2\leq Z_{\beta',\beta}\leq 1-\frac{1}{4}\left(\frac{\xi(\beta')}{\xi(\beta)}\right)^2.
\end{equation}
\end{Prop}

\begin{Prop}[Regularity of the effective random walk below $\beta(\delta)$]
\label{prop: regularity below beta(delta)}
Let $\Creg =3$. There exist $\deltareg>0$ (which can be chosen equal to $2^{-9}$) and $\creg>0$ (which only depends on $d$) such that for every $\beta< \beta(\deltareg)$, the effective random walk at parameter $\beta$ is $(\creg,\Creg)$-regular.
\end{Prop}

\begin{proof}[Proof of Proposition~\textup{\ref{prop: regularity below beta(delta)}}]
The effective random walk $X$ with law $\mathbb P_\beta$ is symmetric by definition.
Let
\begin{equation}
    L:= 2^7 = 128.
    \end{equation}
Let $\delta>0$, to be chosen small enough. In particular, we assume that
$\delta\leq \frac{1}{4L}$.
Fix $\beta< \beta(\delta)$.  Then $\chi(\beta)<\infty$.
Recall that
\begin{equation}
    M_\beta(s)
    =
    \mathbb E_\beta[\exp(s( \mathbf{e}_1\cdot X_1)/\xi(\beta))].
\end{equation}
It suffices to prove that we can choose $\delta,c_1>0$ such
that $M_\beta(c_1)\leq 3$ for every $\beta< \beta(\delta)$.


We define a finite integer $K$ by
\begin{equation}
    K := \min\Big\{k \ge 0 : \text{there exists $\beta_0 < (1-\textstyle\frac{1}{2L})$
    such that $\chi(\beta_0)=L^{-k}\chi(\beta)$}\Big\}.
\end{equation}
%With $\beta_0$ so defined, for $0 \le k \le K$ we define $\beta_k$ by
With $\beta_0$ so defined, for $0 \le k \le K$ we introduce $\beta_k$ by setting
\begin{equation}
    \chi(\beta_k) = L^k \chi(\beta_0).
\end{equation}
In particular, $\beta_K=\beta$.
We will prove by induction on $k$ that $M_{\beta_k}(c_1) \le 3$,
for appropriate $\delta,c_1>0$.

To start the induction, we prove that $M_{\beta_0}(c_1) \le 3$,
with $c_1>0$ to be chosen.
By Proposition~\ref{prop: rough bounds chi xi} and the fact that
(by definition of $\beta(\delta)$)
\begin{equation}
\label{eq:Edelta}
    E(\beta)\leq \delta\leq \frac{1}{4L}\leq \frac{1}{2},
\end{equation}
we have $\xi(\beta_0)\geq \xi(0)=\sigma_J$.
Let $\lambda_0 =\exp[c_1/c_0]$ where $c_0$ is given by \eqref{eq:R_J comparable to xi(0)}.
Then $M_0(c_1\sigma_J/\xi(\beta_0))\leq \exp(c_1\tfrac{\sigma_J}{\xi(\beta_0)}\tfrac{R_J}{\sigma_J})\leq \lambda_0$.
We require that $c_1=c_1(c_0,L,d)>0$ be sufficiently small that
$\lambda_0 \leq 1+\tfrac{1}{4L}$.
Then $\beta_0\lambda_0< (1-\tfrac{1}{2L})(1+\tfrac{1}{4L}) < 1$
and hence, by Corollary~\ref{coro:Miteration},
\begin{equation}
	M_{\beta_0}(c_1)\leq \frac{1}{\chi(\beta_0)}\frac{M_0(c_1\sigma_J/\xi(\beta_0))}{1-\beta_0M_0(c_1\sigma_J/\xi(\beta_0))}
\leq \frac{\lambda_0}{\chi(\beta_0)}\frac{1}{1-\beta_0 \lambda_0}.
\end{equation}
Also, since $\delta\leq \frac{1}{4L}$, \eqref{eq:chi0} implies that
\begin{equation}
	\frac{1}{\chi(\beta_0)}
    =
    \frac{\chi(0)}{\chi(\beta_0)}\leq 1-\beta_0(1-\delta)
    \leq 1-\beta_0\left(1-\frac{1}{4L}\right).
\end{equation}
As a result, we have
\begin{align}
\label{eq:proof reg1}
    M_{\beta_0}(c_1) & \le
    \left(1+\frac{1}{4L}\right)\max_{t\in [0,1-\tfrac{1}{2L}]}
    \frac{1-(1-\tfrac{1}{4L})t}{1-(1+\tfrac{1}{4L})t}
    \nonumber \\
    & =
    \left(1+\frac{1}{4L}\right)
    \frac{1-(1-\tfrac{1}{4L})(1-\tfrac{1}{2L})}{1-(1+\tfrac{1}{4L})(1-\tfrac{1}{2L})}\leq 3.
\end{align}


To advance the induction, we assume that $M_{\beta_k}(c_1) \le C_1$
and prove that the same bound holds when $k$ is replaced by $k+1$.
To abbreviate the notation, we write
\begin{equation}
    Z_k = Z_{\beta_k,\beta_{k+1}}.
\end{equation}
Let
\begin{equation}
    r_k := \frac{\xi(\beta_k)}{\xi(\beta_{k+1})}.
\end{equation}
By Corollary~\ref{coro:Miteration}, if $Z_k M_{\beta_k}(c_1r_k) < 1$ then
\begin{equation}
\label{eq:Msum}
    M_{\beta_{k+1}}(c_1 )
    \le
    \frac{\chi(\beta_k)}{\chi(\beta_{k+1})}
    \frac{M_{\beta_k}(c_1r_k)}{1-Z_k M_{\beta_k}(c_1r_k)}
    .
\end{equation}
To apply this, we must show that $Z_{k} M_{\beta_k}(c_1r_k)<1$.
For this, we will make use of the fact that \eqref{eq:Edelta},
Corollary~\ref{cor: comparison chi xi when E small} and Proposition~\ref{prop:Zbd}
imply the upper bounds
 \begin{align}\label{eq:relation susceptibility-bis}
    r_k^2 \le 2\frac{\chi(\beta_k)}{\chi(\beta_{k+1})}= \frac{2}{L},
\qquad
   Z_{k} &\le  1-\frac{1}{2}\frac{\chi(\beta_k)}{\chi(\beta_{k+1})}=1-\frac{1}{2L}.
    \end{align}
It is elementary that for any $0 \le |u|\le |v|$, we have
\begin{equation}
    \cosh u \le 1 + \frac{u^2}{2} + \frac{u^4}{v^4}\cosh v,
\end{equation}
and hence for any real and symmetric random variable $U$ and
any $0 \le s\le t$,
\begin{equation}\label{eq: general bound laplace transform-bis}
\mathbb E[\exp(sU)] = \mathbb E[\cosh sU]
\le 1+\frac{s^2}{2}\mathbb E[U^2]+\left(\frac st\right)^4\mathbb E[\exp(tU)].
\end{equation}
We apply \eqref{eq: general bound laplace transform-bis}
to $\mathbb E_{\beta_k}$ with $s=c_1r_k$,
$t=c_1$, and $U= (\mathbf{e}_1\cdot X_1)/\xi(\beta_{k})$.
By symmetry, $\mathbb E_{\beta_k}[U^2] =\frac{1}{d}$.
Therefore,  by \eqref{eq:relation susceptibility-bis}
and the induction hypothesis,
\begin{align}
M_{\beta_k}(c_1r_k)
&\le
1+\frac{c_1^2}{2d} r_k^2
+3 r_k^4
\le
1+\frac{c_1^2}{2d}\frac{2}{L}+ 3\left(\frac{2}{L}\right)^{2}
\le 1+\left(c_1^2 + \frac{12}{L}\right)\frac{1}{L} .
\label{eq: proof regularity 2-bis}
\end{align}
We choose $c_1$ small depending on $L$, so that the above gives
\begin{align}\label{eq:regM}
M_{\beta_k}(c_1r_k)
\le 1+  \frac{15}{L} \frac{1}{L}.
\end{align}
Then, by \eqref{eq:relation susceptibility-bis}, and since $\tfrac{15}{L}=\tfrac{15}{128}<1/2$,
\begin{equation}
\label{eq:ZMbd}
    Z_{k}\, M_{\beta_k}(c_1r_k)
    \le
    \left(1-\frac{1}{2L} \right) \left(1 + \frac{15}{L} \frac{1}{L} \right)
    \le
    1 - \left( \frac 12 - \frac{15}{L}\right) \frac 1L
    <1
    .
\end{equation}
This allows us to apply \eqref{eq:Msum}, to obtain
\begin{align}
    M_{\beta_{k+1}}(c_1)
	&
    \le
    \frac 1L \frac{1 + \frac{15}{L} \frac{1}{L}}{( \frac 12 - \frac{15}{L}) \frac 1L}
    =
    \frac{2 + \frac{30}{L^2}}{1 - \frac{30}{L}}
	.
\end{align}
By choice of $L$, the right-hand side is less than
$3$.
This concludes the induction and the proof.

Something important just happened:
we gain by $Z \le 1 -\frac{1}{2L}$ and we lose by $1+\frac{15}{L^2}$ in \eqref{eq: proof regularity 2-bis}
when $c_1$ is small
enough.
When we take the product of these two effects, we still gain.
The potentially bad effect of the constant $15$ is overcome by a choice of large $L$ .



To summarise the two cases, we have obtained that for $L=128$ and $\delta\leq \tfrac{1}{4L}$,  for every $\beta< \beta(\delta)$ we have
\begin{equation}
    M_\beta(c_1)
    \leq 3.
\end{equation}
We now set $\deltareg:=\frac{1}{4L} = 2^{-9}$ and $(\creg,\Creg)=(c_1,3)$.
This completes the proof.
\end{proof}

\end{document}
